%  IKBT User Manual -- using what IKBT produces.
%
%  SOURCE.  The PDF is IKdocs/USER_MANUAL.pdf, built from the repo root with
%
%     pdflatex -output-directory=/tmp LaTex_src/USER_MANUAL.tex   (run twice, for the contents)
%     cp /tmp/USER_MANUAL.pdf IKdocs/
%
%  Every generated report's last section, "Files Produced for <Robot>"
%  (output_latex.generated_files_section()), points here:  that section says
%  WHICH files a robot got, this manual says what to do with any of them.

\documentclass[11pt]{article}
\usepackage[T1]{fontenc}
\usepackage{lmodern}
\usepackage[margin=1in]{geometry}
\usepackage{tabularx}
\usepackage{booktabs}
\usepackage{listings}
\usepackage{enumitem}
\usepackage{xcolor}
\usepackage[obeyspaces,spaces]{url}   % a long command may break at a space
\usepackage[colorlinks=true, linkcolor=blue!50!black, urlcolor=blue!50!black]{hyperref}

\lstset{basicstyle=\ttfamily\small, columns=fullflexible, keepspaces=true,
        upquote=true, frame=single, rulecolor=\color{black!25},
        xleftmargin=1em, xrightmargin=1em, aboveskip=0.8em, belowskip=0.8em}

%  \code{} for identifiers and commands, \file{} for paths.  Both are url-style:
%  _ < > print as typed, and a long one may break at a space, '.', '/' or '_'.
\DeclareUrlCommand\code{\urlstyle{tt}}
\DeclareUrlCommand\file{\urlstyle{tt}}
\newcolumntype{L}{>{\raggedright\arraybackslash}X}

\setlist{itemsep=2pt}
\setlength{\parindent}{0pt}
\setlength{\parskip}{0.6em}

\title{IKBT User Manual\\[0.3em]\large Using what IKBT produces}
\author{}
\date{}

\begin{document}
\maketitle
\raggedright      % long inline commands wrap instead of overrunning

\textbf{Who this is for:}  you have a robot, you have run IKBT on it, and you want
to \emph{use} the result: read the report, call the inverse kinematics from your own program,
 validate the solution, etc.  For knowledge of how IKBT works inside, see the JAIR paper
(\file{IKdocs/IKBT_JAIR_2019.pdf}) and \file{CLAUDE.md}.

\textbf{Which files did my robot get?}  The last section of your robot's report,
\emph{Files Produced for \textless Robot\textgreater}, lists each file
generated by IKBT for your robot.  This manual explains, for every robot, what to do with them.

\subsection*{Requirements:}  Python 3,
\code{sympy}, \code{numpy};  a LaTeX distribution to compile the report;
\code{g++} (C++17) to build the C++.
\tableofcontents

%======================================================================
\section{Running IKBT}

The first step is defining your robot by name in \file{ikbtfunctions/ik_robots.py}
(\code{robot_params()}):  its DH table, its joint variables, and numeric values
(\code{pvals}) for its link lengths.
To add your robot,

\begin{enumerate}
  \item Give it a name which does not appear in the name list ({\tt ikbtfunctions/ik_robots.py: ROBOT_LIST}).
  \item Add the new name to {\tt ROBOT_LIST}.
  \item Copy one of the existing robot definitions (e.g. {\tt if(name == 'KinovaLite'):}),
  and modify it to your specific parameters.
  \item The {\tt vv[]} vector should read {\tt 1} for a rotary joint and {\tt 0} for a prismatic joint.
  \item Your DH table must have 6 rows.  If your robot as $\leq 6$ DOF, add rows of all zeros.
\end{enumerate}

\textbf{Give every parameter a numeric value
in \code{pvals}}. % ;  See Section~\ref{sec:wrong} for what happens if you don't.
Everything runs \textbf{from the IKBT root directory}:

\begin{lstlisting}
python3 ikSolver.py Puma     # solve one robot (in this case ``Puma'')
python3 fkOnly.py Puma       # forward kinematics and Jacobian only, no IK
\end{lstlisting}


A solve takes from seconds (Wrist) to tens of minutes.


%======================================================================
\section{Words and conventions}

These hold for every generated file, in both languages.

\begin{itemize}
\item \textbf{end-effector-pose $T$} ---   a $4\times4$
  homogeneous transform of frame 6 (the last DH frame) in the base frame
  specifying position and orientation of the end effector.
\item \textbf{joint-pose $q$} --- a joint vector, one value per joint, in
  \textbf{DH chain order}.  The generated code names that order in
  \code{JOINT_NAMES}.
\item \textbf{Units.}  Revolute joints are in \textbf{radians}.  Prismatic joints
  and all positions are in the \textbf{length units of the robot's DH table}
  (units are implicit in the numbers used in \code{pvals} ).
\item \textbf{Posture.}  (same as joint-pose)
Most arms reach one end-effector-pose with several
  joint-poses --- elbow up or down, wrist flipped.  IKBT returns \textbf{all} of
  them it can find.  Choosing one is your job, because it depends on things IKBT
  does not know:  joint limits, obstacles, where the arm is now.
\item \textbf{No joint limits}  Returned angles are whatever
  the equations give (for example in $(-\pi, \pi]$).  Apply your joint limits
\item \textbf{Sum-of-angles intermediates} (\code{th_23} and the like) are
  computed internally and \textbf{not} returned;  they are listed in
  \code{AUX_NAMES}.
\end{itemize}

%======================================================================
\section{Which method did my robot get?}\label{sec:method}

As of Version 3.0,
IKBT tries three solution methods in order and stops at the first that works.
The first successful solution method decides which function you
call.

\begin{description}
\item[Symbolic.]  A complete closed form:  an explicit equation for every joint.\\
  Python:  \code{ikin_<Robot>(T)} in \file{IK_equations<Robot>.py}.
  C++:  \code{ikin(T)}.
\item[One-variable.]  No full closed form, but once ONE variable is assumed
  known, the rest is closed form;  a 1-D search finds the right value(s) of that
  variable.  Exact --- nothing about the arm is approximated.\\
  Python:  \code{solve_<Robot>(T)} in \file{IK_onevar<Robot>.py}.
  C++:  \code{solve(T)}.
\item[Hybrid.]  When the first two fail, IKBT creates and solves a slightly \emph{simplified}
  arm in closed form, and uses its solutions as starting points for a 6d numerical
  solution of your actual arm.\\
  Python:  \code{ikin_<Robot>_approx(T)}, then
  \code{refine_seed_<Robot>(T, q_seed)}, in \file{IK_hybrid_<Robot>.py}.
  C++:  \code{ikin_approx(T)}, then \code{refine_seed(T, q_seed)}.

 \item[Fallback]  If the Hybrid method produces an unsatisfactory result, there are many
 full joint space numerical methods out there (ROS Moveit, robotics-toolbox, etc.).  But we do not
 choose one.
\end{description}

To find out which solution method worked on your robot, any of these works:
\begin{itemize}
\item the report's Introduction says ``closed form'', ``ONE-VARIABLE'' or
  ``HYBRID'';
\item the Python file names (above) differ by method;
\item the C++ file says it:
  \code{grep SOLUTION_PATH CodeGen/Cpp/<Robot>CppCode/<Robot>.cpp}.
\end{itemize}

\textbf{Robot names in function names.}  File names use the robot's own name
(\file{IK_onevarC-Arm.py});  The C++ namespace rules require modification of robot
names to eliminate non-identifier characters such as ``-''.  They are
turned into \code{_} (e.g. \code{C-Arm} $\to$ \code{solve_C_Arm},
\code{ikbt::C_Arm}).

%======================================================================
\section{The files, at a glance}

All paths are relative to the IKBT root.  Everything below except
\file{Cpp_src/} is \textbf{generated}:  the next solve of that robot overwrites
it, so edit a copy, never the original.

\begin{center}\small
\begin{tabularx}{\textwidth}{@{}p{0.42\textwidth}Ll@{}}
\toprule
file & what & method \\
\midrule
\file{LaTex/ik_solution_<R>.tex} & the report (Section~\ref{sec:report} of this document) & all \\
\file{CodeGen/Python/IK_equations<R>.py} & closed-form IK & symbolic \\
\file{CodeGen/Python/IK_onevar<R>.py} & the 1-D search --- \textbf{entry point} & one-variable \\
\file{CodeGen/Python/IK_conditional<R>.py} & closed form \emph{given} the assumed variable & one-variable \\
\file{CodeGen/Python/IK_hybrid_<R>.py} & approximate, then correct --- \textbf{entry point} & hybrid \\
\file{CodeGen/Python/IK_equations<D>.py} & closed form of the simplified arm \code{<D>} & hybrid \\
\file{CodeGen/Python/FK_numeric<D>.py} & FK of the simplified arm & hybrid \\
\file{CodeGen/Python/FK_numeric<R>.py} & FK and Jacobian of your robot & all \\
\file{CodeGen/Cpp/<R>CppCode/<R>.cpp} & \textbf{all} the C++ for the robot:  FK, Jacobian, and the IK & all \\
\file{Cpp_src/} & shared C++ headers and a bundled Eigen (source, not generated) & all \\
\file{graphs/<R>_graph.txt} & the solution graph as data & all \\
\file{fk_eqns/<R>_pickle.p} & FK cache;  safe to delete & all \\
\bottomrule
\end{tabularx}
\end{center}

On the hybrid path \code{<D>} is the simplified (``derived'') arm, named for what
was changed --- Raven-II's is \code{Raven-II_a_5_0} (DH parameter \code{a_5} set
to 0).  Its graph file and pickle carry \code{<D>} too.

%======================================================================
\section{Python}

Each generated module requires only numpy  (no IKBT, or sympy imports required) so
you can copy \file{CodeGen/Python/} anywhere.  When using the  one-variable or hybrid
methods, the entry point loads its sibling modules (e.g. {\tt IK_onevar<R>.py})
\textbf{by path}, so keep the
robot's files together in one directory.

Robot names such as \code{C-Arm} are not legal module names, so the examples load
modules by path;  for a name like \code{Puma}, a plain
\code{import IK_equationsPuma} with \file{CodeGen/Python} on \code{sys.path} also
works.

\begin{lstlisting}
import importlib.util, os
import numpy as np

def load(path):
    name = os.path.basename(path)[:-3].replace('-', '_')
    spec = importlib.util.spec_from_file_location(name, path)
    mod = importlib.util.module_from_spec(spec)
    spec.loader.exec_module(mod)
    return mod
\end{lstlisting}

Every module also runs as a script
(\code{python3 CodeGen/Python/IK_equationsPuma.py}) as a self-test.

\subsection{Symbolic}

\begin{lstlisting}
ik = load('CodeGen/Python/IK_equationsPuma.py')
fk = load('CodeGen/Python/FK_numericPuma.py')

T = fk.fk_Puma([0.3, -0.5, 0.4, 0.2, 0.7, -0.1])  # a reachable end-effector-pose
sols = ik.ikin_Puma(T)       # list of joint-poses (JOINT_NAMES order), or False
if sols is False:
    print('unreachable')
else:
    for q in sols:
        print(np.abs(fk.fk_Puma(q) - T).max())  # ~1e-16: every one reaches T
\end{lstlisting}

\begin{itemize}
\item Returns \textbf{one row per posture that exists at $T$}, so the count varies
  with the end-effector-pose;  Puma gives up to 8.  \code{False} means none exist.
\item \code{ikin_Puma_labeled(T)} returns the same rows as dicts keyed by joint
  name.
\item Link lengths are module-level constants (\code{d_1 = 0.6}, \ldots), visible
  at the top of the file.
\end{itemize}

\subsection{One-variable}

\begin{lstlisting}
ov = load('CodeGen/Python/IK_onevarC-Arm.py')
for s in ov.solve_C_Arm(T):
    s['q']            # the joint-pose, JOINT_NAMES order, ready for FK
    s['error']        # ||dp|| + W_ROT*theta at q; accepted only if <= ACCEPT_TOL
    s['known_value']  # the value of KNOWN_VARIABLE it was found at
    s['branch']       # which branch of the closed form it came from
    s['n_samples']    # the search resolution it settled at
\end{lstlisting}

\begin{itemize}
\item The answers are \textbf{exact} --- the equations are your real arm's.
  Typical cost:  about 0.1\,s per call.
\item \code{[]} means no solution was \textbf{found}, which is not quite the same
  as unreachable.
\item If \code{n_samples == MAX_SAMPLES}, the search was still finding new
  solutions when it stopped;  every returned one is still correct, but some may
  be missing.  Raise \code{max_samples} for that end-effector-pose:
  \code{solve_C_Arm(T, max_samples=16384)}.
\item \code{ikin_<R>_given(T, value)}, in \file{IK_conditional<R>.py}, is the
  closed form at one assumed value.  It is exact \textbf{only} at a correct
  value, so do not use it directly to command the arm;  it is exposed for
  plotting and diagnosis (\code{sweep_<R>(T)} gives the raw error curves).
\item Running \file{IK_onevar<R>.py} as a script picks a \textbf{random}
  end-effector-pose and prints the random seed;  pass it back
  (\code{python3 CodeGen/Python/IK_onevarC-Arm.py 12345}) to repeat that case.
\end{itemize}

\subsection{Hybrid --- two calls, on purpose}

\begin{lstlisting}
hy = load('CodeGen/Python/IK_hybrid_Raven-II.py')

approx = hy.ikin_Raven_II_approx(T)       # one approximate joint-pose per posture
q, err, iters = hy.refine_seed_Raven_II(T, approx[0])  # correct ONE
if err <= 1e-9:
    ...                                      # q puts the REAL arm at T

for r in hy.refine_all_Raven_II(T):          # correct every one
    r['index'], r['q'], r['error'], r['converged']
\end{lstlisting}

\begin{itemize}
\item \textbf{Solutions of the  approximate arms are not answers.}
They are good starting points for the numerical method however.
Only a corrected joint-pose, solved by the numerical solver
  reaches $T$.
\item \textbf{\code{refine_seed} takes a joint-pose, not an index}, because
  choosing where to start \emph{is} choosing the posture --- the numerical solver stays
  near its starting point.
\item \textbf{Test the error.}  \code{converged} can legitimately be false:  the
  real arm may not reach $T$ in that posture, or at all.
\item \code{refine_all} is for learning,   which postures the real arm can
  actually reach.
\item The module's header states how the approximate model was derived
(\code{DH_CHANGES}) and its
  cost.
\end{itemize}

\subsection{Forward kinematics and Jacobian (every method)}

\begin{lstlisting}
fk = load('CodeGen/Python/FK_numericPuma.py')
T  = fk.fk_Puma(q)                 # 4x4 numpy array
J6 = fk.jacobian_Puma(q)           # 6 x 6, EXPRESSED IN FRAME 6
J0 = fk.jacobian_base(T, J6)       # the same, rotated into the base frame
\end{lstlisting}

Round-tripping each IK answer through \code{fk_<R>()} is the simplest way to
check it.  On the symbolic path this module is written only when every parameter
has a numeric value.

%======================================================================
\section{C++}

\textbf{\file{IKdocs/CPP_API.md} is the full reference} --- every type and
function, and a compiled-and-run program for each method.  The short version:

\begin{itemize}
\item \textbf{One file per robot}, \file{CodeGen/Cpp/<R>CppCode/<R>.cpp},
  everything in \code{namespace ikbt::<R>}:  \code{fk()}, \code{jacobian()}, and
  the IK (\code{ikin()}, \code{solve()}, or \code{ikin_approx()} +
  \code{refine_seed()}, as in Section~\ref{sec:method}).  Use it by including it:
\begin{lstlisting}
#include "CodeGen/Cpp/PumaCppCode/Puma.cpp"
// T is an Eigen::Matrix4d
std::vector<Eigen::VectorXd> sols = ikbt::Puma::ikin(T);
\end{lstlisting}
\item \textbf{Build:}  \code{g++ -std=c++17 -O2 myapp.cpp -o myapp -lm}.  No
  \code{-I}, nothing to install:  Eigen is bundled in
  \file{Cpp_src/third_party/eigen3/}, and if your build already has Eigen (ROS,
  CMake) the generated code uses yours.
\item \textbf{Self-test:}
\begin{lstlisting}
g++ -std=c++17 -O2 -DIKBT_MAIN CodeGen/Cpp/PumaCppCode/Puma.cpp -o ik && ./ik
\end{lstlisting}
\item \textbf{Keep the directory where it is relative to \file{Cpp_src/}.}  The
  robot file reaches the shared headers by a relative \texttt{\#include}
  (\file{../../../Cpp_src/...}).  To build elsewhere, copy \file{Cpp_src/} along
  with it at the same relative depth.
\item Unreachable is an \textbf{empty vector}, not \code{false}.  The search
  parameters, joint-name arrays and the two hybrid calls mirror the Python one
  for one.
\end{itemize}

%======================================================================
\section{The LaTeX report}\label{sec:report}

\file{LaTex/ik_solution_<R>.tex} is a self-contained document.  To make the PDF:

\begin{lstlisting}
cd LaTex && pdflatex ik_solution_Puma.tex
\end{lstlisting}

It contains the DH parameters, forward kinematics, the unknowns in solution
order, every solution equation (long expressions are shortened by naming pieces
$K_1$, $K_2$, \ldots\ defined just above them), the solution graph (which
variable depends on which), the table of complete solutions, the Jacobian, and
--- on the one-variable and hybrid paths --- sections saying what was assumed or
simplified and how the numerical step works.  The last section lists the files
produced.

The equations are the same ones the generated code evaluates, so the report is
the way to read or cite the solution by hand.  \code{pdflatex} leaves
\file{.aux}, \file{.log} and \file{.out} files beside the PDF;  they can be
deleted.

%======================================================================
\section{Checking the outputs}

Each command runs from the IKBT root.  \code{--keep} checks the files already
written instead of solving the robot again.

\begin{center}\small
\begin{tabularx}{\textwidth}{@{}p{0.40\textwidth}L@{}}
\toprule
command & asks \\
\midrule
\code{python3 -m scripts.numerical_closed_loop_sol_check <R> --keep}
  & Is the generated \textbf{Python} right?  For random reachable
    end-effector-poses, every returned joint-pose must reproduce it
    (\textbf{sound}), and the joint-pose it came from must be among them
    (\textbf{complete}).  Works for all three methods. \\[3pt]
\code{python3 -m scripts.cpp_closed_loop_check <R> --keep}
  & Is the generated \textbf{C++} right, and does it agree with the Python,
    branch for branch? \\[3pt]
\code{python3 -m scripts.check_solution_sets <R>}
  & Is IKBT's \textbf{symbolic} solution set itself right?  Every version must
    reproduce random reachable end-effector-poses, and no two versions may be
    identical. \\
\bottomrule
\end{tabularx}
\end{center}

\textbf{Why \code{python3 -m}?}  These scripts import IKBT's own packages
(\code{ikbtfunctions}, \code{ikbtbasics}, \ldots).  Python puts the
\emph{directory of the script it runs} on its import path, so
\code{python3 scripts/numerical_closed_loop_sol_check.py} looks for those
packages inside \file{scripts/} and stops with
\code{ModuleNotFoundError: No module named 'ikbtfunctions'}.
\code{python3 -m scripts.<name>}, run from the IKBT root, puts the root on the
path instead, and everything is found.

\file{IKdocs/TESTING.md} describes these and the rest of the test suite in
detail.

%======================================================================
\section{Other outputs}

These are made only when you run them;  none of them solves the robot again.

\begin{center}\small
\begin{tabularx}{\textwidth}{@{}p{0.32\textwidth}L@{}}
\toprule
how & what \\
\midrule
\code{python3 -m scripts.tube_map <R>}
  & \file{graphs/<R>_tube.svg}:  the solution graph drawn as a subway map ---
    one coloured line per complete solution, one station per version of each
    variable.  Reads only \file{graphs/<R>_graph.txt}, so it is instant.  On
    the hybrid path use the simplified arm's name. \\[3pt]
\code{python3 fkOnly.py <R>}
  & forward kinematics only:  \file{LaTex/fk_equations_<R>.tex},
    \file{CodeGen/Python/FK_equations<R>.py} (a readable dump of the symbolic
    FK), and \file{CodeGen/Cpp/<R>CppCode/FK_numeric<R>.h}.  Use it for a robot
    you need only FK for;  it does not touch IK files. \\[3pt]
\code{python3 -m scripts.bt_diagram}
  & the behavior tree itself, as slides (\file{IKdocs/IKBT_BT_p1..p3.svg}). \\
\bottomrule
\end{tabularx}
\end{center}

%======================================================================
\section{Regenerating and cleaning up}

\begin{itemize}
\item \textbf{Re-run the solve} to regenerate everything for a robot;  it
  overwrites that robot's files.
\item \textbf{\file{CodeGen/cleanCodeGenOutput}} (run from the root) deletes all
  generated Python and C++.
\item \textbf{\file{LaTex_src/cleanLaTexFolder}} (run from the root) deletes all
  generated reports.
\item \textbf{\file{fk_eqns/}} is a cache of forward kinematics;  delete any or
  all of it at any time.  The next solve recomputes it, which costs only time.
  Changing a robot's DH table refreshes its cache automatically.
\item \file{Cpp_src/} and \file{LaTex_src/} are source, not output:  the cleanup
  scripts never touch them, and neither should you.
\end{itemize}

%======================================================================
\section{When something looks wrong}\label{sec:wrong}

\begin{center}\small
\begin{tabularx}{\textwidth}{@{}p{0.36\textwidth}L@{}}
\toprule
you see & it means \\
\midrule
\code{ikin_<R>()} returns \code{False} (Python) or an empty vector (C++)
  & no posture exists at that end-effector-pose.  Check it is in the workspace,
    and in the DH table's units. \\[3pt]
fewer joint-poses than the report's solution table has rows
  & normal:  rows whose \code{acos}/\code{asin}/\code{sqrt} argument leaves its
    domain at this end-effector-pose are postures that do not exist there, and
    are dropped. \\[3pt]
\code{XXXXX} in the C++, and \code{g++} stops on it
  & a parameter has no \code{pvals} entry.  Deliberate:  a missing link length
    is never silently defaulted.  Add it to \file{ik_robots.py} and re-solve, or
    type the value in. \\[3pt]
no \file{FK_numeric<R>.py} after a symbolic solve
  & same cause:  FK needs every parameter to have a value.  The IK is still
    fine. \\[3pt]
hybrid:  \code{converged} false / large error
  & the real arm does not reach $T$ from that starting joint-pose.  Try
    another, or the arm's current joint-pose;  if none converges, $T$ may be
    unreachable. \\[3pt]
one-variable:  \code{n_samples == MAX_SAMPLES}
  & the search had not settled;  raise \code{max_samples}. \\[3pt]
one-variable:  \code{[]} for an end-effector-pose you believe reachable
  & raise \code{n_samples}/\code{max_samples}, and plot \code{sweep_<R>(T)} to
    see the error curves. \\[3pt]
a solve takes very long
  & normal for hard arms (minutes).  \code{python3 ikSolver.py <R> --perf}
    reports sympy timings and names each slow call as it happens. \\
\bottomrule
\end{tabularx}
\end{center}

\end{document}
